% GENERATED FILE. Edit canonical lessons, then run npm run generate:books.
% canonical-source-hash: fnv1a64:0cb52057511d074a
% canonical-lessons: LA-C443-syrtis, LA-C443-fluctus, LA-C443-vadum, LA-C443-fretum, LA-C443-praetervehor

\chapter{Sandbank, Billow, Shallows, Strait, To sail past}
\label{ch:la-sandbank-billow-shallows-strait-to-sail-past}
\hlchaptermodality{443}

\begin{chapteropening}
\textbf{By the end of this chapter:} \emph{I can talk about the dangers of the water: a sandbank, the billows, the shallows, a strait, and sailing past.}

Complete the last lesson of chapter 443: I can talk about the dangers of the water: a sandbank, the billows, the shallows, a strait, and sailing past.
\end{chapteropening}

\section[syrtis, syrtis]{syrtis, syrtis — a sandbank, a shoal}

\label{lesson:LA-C443-syrtis}



\begin{quote}
Before the new one: say the Latin for of the sea, then the Latin for coastal.
\end{quote}

\subsection*{You'll want to know: syrtis, syrtis}
\textbf{syrtis, syrtis} — \textquotedblleft{}a sandbank, a shoal\textquotedblright{}.

\begin{quote}
Gubernātor syrtēs vītat et nāvem in altum dūcit.
\end{quote}

\begin{quote}
\emph{The helmsman avoids the sandbanks and steers the ship into deep water.}
\end{quote}

Third declension, feminine: \emph{syrtis, syrtis f.} Any sandbank, and above all the two dangerous sandy gulfs off the African coast.

\begin{cousinweb}
\emph{syrtis} was borrowed from Greek \emph{syrtis} 'sandbank' (Lewis \& Short), above all the two gulfs off Africa full of shoals. EtymDB lists it without a parent.
\end{cousinweb}

\begin{grammarlens}[title={What you've built}]
One more word for telling, explaining and describing at B1.
\end{grammarlens}

\subsection*{Your turn}
\begin{itemize}
\raggedright
  \item \emph{Say it:} \emph{syrtis, syrtis}
  \item \emph{Say it:} \emph{syrtis, syrtis}, once more
  \item \emph{Say it:} a sentence of your own with \emph{syrtis, syrtis}
  \item \emph{Recall:} say the Latin for of the sea, then the Latin for coastal, then say \emph{syrtis, syrtis} again
\end{itemize}

\begin{tcolorbox}[breakable,colback=teal!4,colframe=teal!35!black,title={Before you move on}]
What does syrtis, syrtis mean? (\textquotedblleft{}A sandbank, a shoal\textquotedblright{}.) Say it once more.
\end{tcolorbox}
\section[flūctus, flūctūs]{flūctus, flūctūs — a billow, a wave; a surge}

\label{lesson:LA-C443-fluctus}



\begin{quote}
Before the new one: say the Latin for coastal, then the Latin for a sandbank.
\end{quote}

\subsection*{You'll want to know: flūctus, flūctūs}
\textbf{flūctus, flūctūs} — \textquotedblleft{}a billow, a wave; a surge\textquotedblright{}.

\begin{quote}
Magnī flūctūs in lītus cadunt.
\end{quote}

\begin{quote}
\emph{Great waves fall on the shore.}
\end{quote}

Fourth declension, masculine: \emph{flūctus, flūctūs m.} A bigger word than \emph{unda}: the wave as it rolls and breaks.

\begin{cousinweb}
\emph{flūctus} is built on \emph{fluō} 'to flow' with the noun suffix \emph{-tus} (EtymDB; Lewis \& Short): a flowing, a surge of water.
\end{cousinweb}

\begin{grammarlens}[title={What you've built}]
One more word for telling, explaining and describing at B1.
\end{grammarlens}

\subsection*{Your turn}
\begin{itemize}
\raggedright
  \item \emph{Say it:} \emph{flūctus, flūctūs}
  \item \emph{Say it:} \emph{flūctus, flūctūs}, once more
  \item \emph{Say it:} a sentence of your own with \emph{flūctus, flūctūs}
  \item \emph{Recall:} say the Latin for coastal, then the Latin for a sandbank, then say \emph{flūctus, flūctūs} again
\end{itemize}

\begin{tcolorbox}[breakable,colback=teal!4,colframe=teal!35!black,title={Before you move on}]
What does flūctus, flūctūs mean? (\textquotedblleft{}A billow, a wave; a surge\textquotedblright{}.) Say it once more.
\end{tcolorbox}
\section[vadum, vadī]{vadum, vadī — a shallow, a shoal; a ford}

\label{lesson:LA-C443-vadum}



\begin{quote}
Before the new one: say the Latin for a sandbank, then the Latin for a billow.
\end{quote}

\subsection*{You'll want to know: vadum, vadī}
\textbf{vadum, vadī} — \textquotedblleft{}a shallow, a shoal; a ford\textquotedblright{}.

\begin{quote}
Equī flūmen per vadum trānseunt.
\end{quote}

\begin{quote}
\emph{The horses cross the river by the ford.}
\end{quote}

Second declension, neuter: \emph{vadum, vadī n.} A place where water is shallow enough to walk through, or to run aground on.

\begin{cousinweb}
\emph{vadum} is inherited from Proto-Italic \emph{waðom} 'shallow, ford', from PIE \emph{wh\textsubscript{2}d\textsuperscript{h}óm} (EtymDB). Lewis \& Short link it with \emph{vādō} 'to go': a place where you can walk through water.
\end{cousinweb}

\begin{grammarlens}[title={What you've built}]
One more word for telling, explaining and describing at B1.
\end{grammarlens}

\subsection*{Your turn}
\begin{itemize}
\raggedright
  \item \emph{Say it:} \emph{vadum, vadī}
  \item \emph{Say it:} \emph{vadum, vadī}, once more
  \item \emph{Say it:} a sentence of your own with \emph{vadum, vadī}
  \item \emph{Recall:} say the Latin for a sandbank, then the Latin for a billow, then say \emph{vadum, vadī} again
\end{itemize}

\begin{tcolorbox}[breakable,colback=teal!4,colframe=teal!35!black,title={Before you move on}]
What does vadum, vadī mean? (\textquotedblleft{}A shallow, a shoal; a ford\textquotedblright{}.) Say it once more.
\end{tcolorbox}
\section[fretum, fretī]{fretum, fretī — a strait, a narrow channel; (in poetry) the sea}

\label{lesson:LA-C443-fretum}



\begin{quote}
Before the new one: say the Latin for a billow, then the Latin for a shallow.
\end{quote}

\subsection*{You'll want to know: fretum, fretī}
\textbf{fretum, fretī} — \textquotedblleft{}a strait, a narrow channel; (in poetry) the sea\textquotedblright{}.

\begin{quote}
Fretum angustum est, sed flūctūs magnī sunt.
\end{quote}

\begin{quote}
\emph{The strait is narrow, but the waves are big.}
\end{quote}

Second declension, neuter: \emph{fretum, fretī n.} A narrow sea between two coasts, where the water runs hard.

\begin{cousinweb}
\emph{fretum} is inherited. Lewis \& Short connect it with a root for restless, seething movement, seen in \emph{ferveō} 'to boil', since water churns in a strait; EtymDB gives no parent, so the link is hedged.
\end{cousinweb}

\begin{grammarlens}[title={What you've built}]
One more word for telling, explaining and describing at B1.
\end{grammarlens}

\subsection*{Your turn}
\begin{itemize}
\raggedright
  \item \emph{Say it:} \emph{fretum, fretī}
  \item \emph{Say it:} \emph{fretum, fretī}, once more
  \item \emph{Say it:} a sentence of your own with \emph{fretum, fretī}
  \item \emph{Recall:} say the Latin for a billow, then the Latin for a shallow, then say \emph{fretum, fretī} again
\end{itemize}

\begin{tcolorbox}[breakable,colback=teal!4,colframe=teal!35!black,title={Before you move on}]
What does fretum, fretī mean? (\textquotedblleft{}A strait, a narrow channel; (in poetry) the sea\textquotedblright{}.) Say it once more.
\end{tcolorbox}
\section[praetervehor, praetervehī]{praetervehor, praetervehī — to sail past, to ride past, to be carried past}

\label{lesson:LA-C443-praetervehor}



\begin{quote}
Before the new one: say the Latin for a shallow, then the Latin for a strait.
\end{quote}

\subsection*{You'll want to know: praetervehor, praetervehī}
\textbf{praetervehor, praetervehī} — \textquotedblleft{}to sail past, to ride past, to be carried past\textquotedblright{}.

\begin{quote}
Nāvis multa oppida praetervehitur.
\end{quote}

\begin{quote}
\emph{The ship sails past many towns.}
\end{quote}

Deponent in form: \emph{praetervehor, praetervehī, praetervectus sum}, with an accusative of the place passed.

\begin{cousinweb}
\emph{praetervehor} is built from \emph{praeter} 'past, beyond' and \emph{vehor}, the passive of \emph{vehō} 'to carry' (Lewis \& Short): to be carried past a place by ship or carriage.
\end{cousinweb}

\begin{grammarlens}[title={What you've built}]
One more word for telling, explaining and describing at B1.
\end{grammarlens}

\subsection*{Your turn}
\begin{itemize}
\raggedright
  \item \emph{Say it:} \emph{praetervehor, praetervehī}
  \item \emph{Say it:} \emph{praetervehor, praetervehī}, once more
  \item \emph{Say it:} a sentence of your own with \emph{praetervehor, praetervehī}
  \item \emph{Recall:} say the Latin for a shallow, then the Latin for a strait, then say \emph{praetervehor, praetervehī} again
\end{itemize}

\begin{tcolorbox}[breakable,colback=teal!4,colframe=teal!35!black,title={Before you move on}]
What does praetervehor, praetervehī mean? (\textquotedblleft{}To sail past, to ride past, to be carried past\textquotedblright{}.) Say it once more.
\end{tcolorbox}
